\section{Evaluation Framework}
\label{sec:7-1}

This chapter estimates the \textbf{business value} of targeting with the final model of Chapter~\ref{ch:modeling}, whose detection quality is reported in Section~\ref{sec:6-quality}. All evaluation uses the \textbf{Scoring Population}: every active customer at every weekly test date (10 weeks, 2011-07-04 -- 2011-09-05, 24,252 rows). The same customer appears in several consecutive weeks, so uncertainty is reported across weeks rather than as if rows were independent. Repeat customers make up 86.7\% of test rows and one-time customers 13.3\%; segment-level error rates are in Section~\ref{sec:7-12}. The historical data contain no campaigns, so suppression is not part of the data; Section~\ref{sec:7-sens} simulates it -- removing customers contacted in the previous $S = 13$ weeks -- as one of three ways to count a weekly campaign.

\textbf{Candidates.} Besides the final model and the three baselines (recency rule, RFM rule, random), this chapter evaluates a value-aware ranking, \textbf{Expected Loss} $=$ calibrated churn probability $\times$ revenue at risk ($V \times$ expected orders, Section~\ref{sec:7-5}). It uses the same fitted model; only the ranking key changes, from ``who is most likely to leave'' to ``whose departure would cost the most''.

\section{Customer Revenue at Risk}
\label{sec:7-5}

Estimate the future revenue a churner's departure would cost, using history up to the snapshot date only. These are estimates of revenue exposure, not observed lost revenue. Three measures are compared throughout:

\begin{center}
\begin{tabular}{@{}p{2.6cm}p{5.6cm}p{6cm}@{}}
    \toprule
    Measure & Definition & Rationale \\
    \midrule
    \texttt{rar\_aov} & One future purchase at the raw average order value so far & The reference measure of earlier versions \\
    \texttt{rar\_v} & One future purchase at the capped value $V$ (Section~\ref{sec:5-5}) & A handful of wholesale-sized accounts cannot dominate \\
    \texttt{rar\_v\_orders} & $V \times \text{clip}(H / \text{mean gap so far},\ 1,\ 13)$ & Losing a frequent buyer costs more orders within $H$ than losing an occasional one; one-time customers count one order \\
    \bottomrule
\end{tabular}
\end{center}

\begin{center}
\begin{tabular}{@{}p{3.6cm}p{1.8cm}p{1.8cm}p{1.8cm}p{2.2cm}@{}}
    \toprule
    Measure & Mean & Median & 75th pct. & Max \\
    \midrule
    \texttt{rar\_aov} & 395.70 & 303.58 & 438.93 & 14,844.77 \\
    \texttt{rar\_v} & 373.69 & 303.58 & 438.93 & 1,956.84 \\
    Expected orders in $H$ & 2.07 & 1.20 & 2.20 & 13.00 \\
    \texttt{rar\_v\_orders} & 832.85 & 448.88 & 844.24 & 25,438.92 \\
    \bottomrule
\end{tabular}
\end{center}

Capping only affects the tail (the median is unchanged), but the order-count factor more than doubles the mean exposure and spreads it widely: a customer buying every week can carry 13 times the exposure of a one-time buyer with the same basket.

\section{Business Impact Scenario}
\label{sec:7-6}

For a target share $K$, per scoring week: customers targeted ($N_K$), true churners captured ($TP_K$), non-churners targeted ($FP_K$), Recall@K, and revenue at risk among targeted churners. Sections~\ref{sec:7-6}--\ref{sec:7-9} keep the \textbf{reference measurement} of earlier versions -- weekly counting, \texttt{rar\_aov} -- so the funnel stays readable; Section~\ref{sec:7-sens} shows how much the answer depends on both choices.

At $K=10\%$, summed across the 10 test weeks (so $N_K$ counts weekly targeting instances, not unique customers -- the same customer can be targeted in more than one week):

\begin{center}
\begin{tabular}{@{}p{2.8cm}p{1.6cm}p{1.6cm}p{1.6cm}p{2.8cm}@{}}
    \toprule
    Candidate & $N_K$ & $TP_K$ & $FP_K$ & Revenue at risk \\
    \midrule
    Final model & 2,427 & 1,831 & 596 & 406,981.23 \\
    Expected Loss & 2,427 & 1,132 & 1,295 & 1,165,360.89 \\
    RFM rule & 2,427 & 1,732 & 695 & 499,916.48 \\
    Recency rule & 2,427 & 1,521 & 906 & 538,397.09 \\
    Random & 2,427 & 986 & 1,441 & 330,372.57 \\
    \bottomrule
\end{tabular}
\end{center}

$N_K$ is identical across candidates by construction (same $K$, same weekly population); $TP_K$ and revenue at risk are not -- the final model captures the most churners, but the two baselines happen to target higher-average-order-value customers, so they show \emph{more} revenue at risk despite capturing fewer churners. Expected Loss pushes this to the extreme: the fewest churners of the four learned/rule-based rankings, but almost three times the final model's revenue at risk. This tension carries through to Sections~\ref{sec:7-9}--\ref{sec:7-10}.

\section{Retention Effectiveness Assumption}
\label{sec:7-7}

Retention success is an assumption, not an observed result. Scenarios: \textbf{Conservative $e = 30\%$}, \textbf{Base $e = 40\%$}, \textbf{Optimistic $e = 50\%$}.
\begin{equation*}
    \text{RevenuePreserved} = \text{RevenueAtRisk(targeted churners)} \times e
\end{equation*}

At $K=10\%$ (Section~\ref{sec:7-6}'s targeted churners), revenue preserved across the three effectiveness scenarios:

\begin{center}
\begin{tabular}{@{}p{2.8cm}p{2.6cm}p{2.6cm}p{2.6cm}@{}}
    \toprule
    Candidate & $e=30\%$ & $e=40\%$ (Base) & $e=50\%$ \\
    \midrule
    Final model & 122,094.37 & 162,792.49 & 203,490.62 \\
    Expected Loss & 349,608.27 & 466,144.36 & 582,680.45 \\
    RFM rule & 149,974.94 & 199,966.59 & 249,958.24 \\
    Recency rule & 161,519.13 & 215,358.84 & 269,198.55 \\
    Random & 99,111.77 & 132,149.03 & 165,186.28 \\
    \bottomrule
\end{tabular}
\end{center}

Because revenue preserved scales linearly with $e$ for a fixed set of targeted churners, the ranking across candidates is unchanged by the effectiveness assumption -- only Section~\ref{sec:7-9}'s cost terms can shift it.

\section{Marketing Cost and the Value of the Model}
\label{sec:7-8}

If contact cost per customer is roughly constant, targeting $K\%$ instead of everyone reduces reach cost by $(1 - K)$. This saving comes from \textbf{targeting itself}, not from the model -- any ranking at the same $K$ saves the same amount.

The model's contribution is measured by comparing, \textbf{at the same $K$}, the churners captured, revenue preserved and net impact of:
\begin{enumerate}
    \item the final model,
    \item the recency-rule baseline, and
    \item random targeting.
\end{enumerate}

At $K=10\%$, $e=40\%$ (Base), contact cost is identical across candidates by construction ($N_K \times c_{\text{contact}}$, same $N_K$):

\begin{center}
\begin{tabular}{@{}p{2.8cm}p{3cm}p{3cm}p{2.6cm}@{}}
    \toprule
    Candidate & Customers targeted & Churners captured & Contact cost \\
    \midrule
    Final model & 2,427 & 1,831 & 1,213.50 \\
    Expected Loss & 2,427 & 1,132 & 1,213.50 \\
    RFM rule & 2,427 & 1,732 & 1,213.50 \\
    Recency rule & 2,427 & 1,521 & 1,213.50 \\
    Random & 2,427 & 986 & 1,213.50 \\
    \bottomrule
\end{tabular}
\end{center}

The reach-cost saving from targeting 10\% instead of everyone is the same 1,213.50 regardless of which candidate ranks customers -- confirming it comes from \emph{targeting}, not the model. The model's own contribution shows up only in churners captured (1,831 vs.\ 986 for random, almost double) -- whether that translates into more net business value depends on revenue at risk and incentive cost too (Section~\ref{sec:7-9}).

\section{Net Business Impact}
\label{sec:7-9}

Because preserved revenue is not profit, a margin assumption converts it before subtracting costs:
\begin{align*}
    \text{GrossProfitPreserved} &= \text{RevenuePreserved} \times m \\
    \text{ContactCost} &= N_K \times c_{\text{contact}} \\
    \text{IncentiveCost} &= (TP_K \times e + FP_K \times r) \times c_{\text{inc}} \\
    \text{NetBusinessImpact} &= \text{GrossProfitPreserved} - \text{ContactCost} - \text{IncentiveCost}
\end{align*}
Here $r$ is the assumed share of targeted non-churners who redeem an incentive although they would have purchased anyway (cannibalization). Optionally report $\text{GrossProfitPreserved} \div (\text{ContactCost} + \text{IncentiveCost})$.

At $K=10\%$, $e=40\%$ (Base), $m=35\%$, $c_{\text{contact}}=\pounds0.50$, $c_{\text{inc}}=\pounds5.00$, $r=20\%$:

\begin{center}
\footnotesize
\setlength{\tabcolsep}{3pt}
\begin{tabular}{@{}p{2.2cm}p{1.3cm}p{1.3cm}p{1.6cm}p{1.9cm}p{1.6cm}p{1.6cm}p{1.9cm}@{}}
    \toprule
    Candidate & $TP_K$ & $FP_K$ & Revenue preserved & Gross profit preserved & Contact cost & Incentive cost & Net impact \\
    \midrule
    Final model & 1,831 & 596 & 162,792.49 & 56,977.37 & 1,213.50 & 4,258.00 & 51,505.87 \\
    Expected Loss & 1,132 & 1,295 & 466,144.36 & 163,150.52 & 1,213.50 & 3,559.00 & \textbf{158,378.02} \\
    RFM rule & 1,732 & 695 & 199,966.59 & 69,988.31 & 1,213.50 & 4,159.00 & 64,615.81 \\
    Recency rule & 1,521 & 906 & 215,358.84 & 75,375.59 & 1,213.50 & 3,948.00 & 70,214.09 \\
    Random & 986 & 1,441 & 132,149.03 & 46,252.16 & 1,213.50 & 3,413.00 & 41,625.66 \\
    \bottomrule
\end{tabular}
\end{center}

\textbf{Counter-intuitive result.} At this operating point, the \emph{final model's net business impact is below both rule-based baselines} -- even below the simple recency rule -- despite capturing the most churners and leading every ranking metric of Section~\ref{sec:6-quality}. The reason is visible in Section~\ref{sec:7-6}: revenue at risk is driven by \emph{which} customers are targeted (their average order value), not by how many churners are caught. The recency and RFM rules happen to target a higher-order-value subset of customers, so they preserve more revenue per targeted churner even while catching fewer of them. This is not a bug in the model -- it optimizes churn-probability ranking, not revenue-weighted ranking -- but it means \textbf{the best classifier is not automatically the best business decision at a fixed $K$}. Ranking the same model's output by value instead (Expected Loss) triples net impact under this measurement, to \pounds158,378. That figure is inflated by the reference measurement itself, however: Expected Loss reaches only 214 distinct churners and counts each about five times, once per week they stay in its top 10\% -- which is why Section~\ref{sec:7-sens} changes how the campaign is counted.

\section{Measurement Sensitivity: Value and Counting}
\label{sec:7-sens}

Two choices behind Sections~\ref{sec:7-6}--\ref{sec:7-9} are open to question: the revenue measure (Section~\ref{sec:7-5}) and how a weekly campaign is counted. Three counting schemes are compared:

\begin{itemize}
    \item \textbf{Weekly} (reference): the top $K\%$ of all active customers every week. A customer who stays at the top is counted -- and paid for -- every week, although they can only be retained once per churn episode.
    \item \textbf{Unique}: the same weekly selection, but each customer counts once, at their first targeting. This removes double counting but is \emph{not} budget-equal: a ranking that keeps re-selecting the same customers ends up spending on fewer people.
    \item \textbf{Suppressed}: every week the same budget $N_K$ goes to the top-ranked customers \emph{not contacted in the previous 13 weeks} ($\approx H$), the suppression rule of Section~\ref{sec:8-5}. This is both budget-equal and free of double counting -- the operationally realistic scheme.
\end{itemize}

Net impact at $K=10\%$, $e=40\%$, for every combination (``Churners'' is the number of distinct churners reached):

\begin{center}
\footnotesize
\setlength{\tabcolsep}{4pt}
\begin{tabular}{@{}p{1.6cm}p{2.2cm}p{1.3cm}p{1.1cm}p{1.4cm}p{1.8cm}p{1.8cm}p{2cm}@{}}
    \toprule
    Counting & Candidate & Targeted & $TP_K$ & Churners & \texttt{rar\_aov} & \texttt{rar\_v} & \texttt{rar\_v\_orders} \\
    \midrule
    Weekly & Final model & 2,427 & 1,831 & 425 & 51,506 & 51,506 & 64,698 \\
     & Expected Loss & 2,427 & 1,132 & 214 & \textbf{158,378} & \textbf{112,113} & \textbf{270,653} \\
     & Recency rule & 2,427 & 1,521 & 542 & 70,214 & 63,767 & 94,395 \\
     & RFM rule & 2,427 & 1,732 & 531 & 64,616 & 64,616 & 84,222 \\
     & Random & 2,427 & 986 & 719 & 41,626 & 41,329 & 72,988 \\
    \midrule
    Unique & Final model & 577 & 425 & 425 & 13,679 & 13,679 & 17,249 \\
     & Expected Loss & 431 & 199 & 199 & 24,192 & 19,090 & 44,647 \\
     & Recency rule & 899 & 542 & 542 & 26,072 & 23,234 & 33,790 \\
     & RFM rule & 771 & 531 & 531 & 20,134 & 20,134 & 25,920 \\
     & Random & 1,635 & 679 & 679 & \textbf{28,731} & \textbf{28,625} & \textbf{47,586} \\
    \midrule
    Suppressed & Final model & 2,427 & 1,212 & 1,212 & \textbf{54,830} & \textbf{49,787} & 77,305 \\
     & Expected Loss & 2,427 & 1,000 & 1,000 & 54,513 & 49,411 & \textbf{81,337} \\
     & Recency rule & 2,427 & 1,152 & 1,152 & 52,865 & 48,144 & 75,114 \\
     & RFM rule & 2,427 & 1,220 & 1,220 & 51,350 & 49,700 & 70,273 \\
     & Random & 2,427 & 986 & 986 & 47,364 & 42,321 & 71,258 \\
    \bottomrule
\end{tabular}
\end{center}

Four observations:
\begin{enumerate}
    \item \textbf{Weekly counting inflates every total, unequally.} The final model reaches only 425 distinct churners behind its 1,831 $TP_K$ (4.3 weeks each), Expected Loss 214 behind 1,132 (5.3 each), random 719 behind 986. Rankings that keep the same customers at the top are rewarded for it.
    \item \textbf{Capping matters only for the recency rule.} Four targeted wholesale-sized rows carry \pounds54K of its raw revenue at risk; with $V$ its weekly net impact falls from \pounds70,214 to \pounds63,767, below the RFM rule. The final model and the RFM rule never reach those accounts.
    \item \textbf{Unique counting mostly measures reach.} Because it is not budget-equal, random -- which spreads its selections over 1,635 customers -- comes out on top; the ranking that re-selects the most (the final model, 577 customers) comes out last.
    \item \textbf{Under suppressed counting at $K=10\%$, the candidates converge} (\pounds70--81K under \texttt{rar\_v\_orders}): with a 13-week suppression over a 10-week test, a 10\% weekly budget contacts 83\% of all 2,907 test customers at least once, so every ranking ends up reaching almost the same people, only in a different order. Section~\ref{sec:7-10} therefore looks at smaller weekly shares.
\end{enumerate}

\section{Scenario Analysis}
\label{sec:7-10}

Net impact under the realistic scheme -- suppressed counting -- for weekly target shares below the saturation point, all three revenue measures, $e=40\%$ and the cost assumptions of Section~\ref{sec:7-9}. ``Reach'' is the share of the 2,907 test customers contacted at least once; the last column compares each candidate with the recency rule under \texttt{rar\_v\_orders}:

\begin{center}
\footnotesize
\setlength{\tabcolsep}{4pt}
\begin{tabular}{@{}p{0.8cm}p{2.3cm}p{1cm}p{1.1cm}p{1.8cm}p{1.8cm}p{2.1cm}p{2.1cm}@{}}
    \toprule
    $K$ & Candidate & Reach & $TP_K$ & \texttt{rar\_aov} & \texttt{rar\_v} & \texttt{rar\_v\_orders} & vs.\ recency \\
    \midrule
    1\% & Final model & 8\% & 185 & 3,926 & 3,926 & 4,899 & $-$4,536 \\
     & Expected Loss & 8\% & 113 & \textbf{17,660} & \textbf{12,558} & \textbf{35,258} & $+$25,823 \\
     & Recency rule & 8\% & 166 & 6,157 & 6,157 & 9,435 & -- \\
     & RFM rule & 8\% & 188 & 6,501 & 6,501 & 6,969 & $-$2,465 \\
     & Random & 8\% & 100 & 4,882 & 4,690 & 7,462 & $-$1,972 \\
    \midrule
    2\% & Final model & 17\% & 363 & 10,776 & 10,776 & 12,782 & $-$7,337 \\
     & Expected Loss & 17\% & 225 & \textbf{25,704} & \textbf{20,574} & \textbf{46,366} & $+$26,247 \\
     & Recency rule & 17\% & 315 & 15,057 & 13,253 & 20,120 & -- \\
     & RFM rule & 17\% & 362 & 13,581 & 13,581 & 17,963 & $-$2,157 \\
     & Random & 17\% & 189 & 7,964 & 7,964 & 15,740 & $-$4,380 \\
    \midrule
    5\% & Final model & 42\% & 766 & 28,866 & 27,832 & 37,031 & $-$4,164 \\
     & Expected Loss & 42\% & 517 & \textbf{40,594} & \textbf{35,493} & \textbf{65,354} & $+$24,159 \\
     & Recency rule & 42\% & 688 & 31,175 & 28,278 & 41,195 & -- \\
     & RFM rule & 42\% & 742 & 29,214 & 28,181 & 38,051 & $-$3,144 \\
     & Random & 42\% & 496 & 21,452 & 21,347 & 38,305 & $-$2,890 \\
    \midrule
    10\% & Final model & 83\% & 1,212 & 54,830 & 49,787 & 77,305 & $+$2,191 \\
     & Expected Loss & 83\% & 1,000 & 54,513 & 49,411 & 81,337 & $+$6,224 \\
     & Recency rule & 83\% & 1,152 & 52,865 & 48,144 & 75,114 & -- \\
     & RFM rule & 83\% & 1,220 & 51,350 & 49,700 & 70,273 & $-$4,841 \\
     & Random & 83\% & 986 & 47,364 & 42,321 & 71,258 & $-$3,855 \\
    \bottomrule
\end{tabular}
\end{center}

Below saturation, the picture is consistent across all three revenue measures. \textbf{Expected Loss is the best ranking at every $K$ from 1\% to 5\%}, by a wide margin -- at $K=2\%$, \pounds46,366 against \pounds20,120 for the recency rule and \pounds12,782 for the final model under \texttt{rar\_v\_orders}, and still ahead under the raw order value it was not built on (\pounds25,704 vs.\ \pounds15,057 and \pounds10,776). The final model, ranking by churn probability alone, is the \emph{weakest} ranking at these budgets under \texttt{rar\_v\_orders} -- below even random targeting: it spends a scarce budget on customers who are almost certain to leave but would cost little (mostly one-time buyers, Section~\ref{sec:7-12}). Expected Loss catches fewer churners ($TP_K$ 225 vs.\ 363 at $K=2\%$) -- its targets churn less often -- but the cost of contacting a non-churner (\pounds0.50 plus a 20\% chance of a \pounds5 redemption) is small next to the value of keeping one high-value churner. One caveat: Expected Loss is built on \texttt{rar\_v\_orders}, so its lead under that measure is partly by construction; its lead under \texttt{rar\_aov} and \texttt{rar\_v} is the independent evidence.

\section{Model Comparison}
\label{sec:7-11}

Compare candidates (including baselines) on detection quality and business-impact scenarios. Do not define a single "best" model from classification metrics alone when the objective depends on ranking and cost constraints.

\begin{center}
\begin{tabular}{@{}p{3cm}p{1.6cm}p{4.2cm}p{4.2cm}@{}}
    \toprule
    Candidate & AUC & Net impact, reference (weekly, \texttt{rar\_aov}, $K$=10\%) & Net impact, realistic (suppressed, \texttt{rar\_v\_orders}, $K$=5\%) \\
    \midrule
    Final model & \textbf{0.759} & 51,506 & 37,031 \\
    Expected Loss & 0.533 & \textbf{158,378} & \textbf{65,354} \\
    RFM rule & 0.728 & 64,616 & 38,051 \\
    Recency rule & 0.638 & 70,214 & 41,195 \\
    Random & 0.506 & 41,626 & 38,305 \\
    \bottomrule
\end{tabular}
\end{center}

No candidate wins on every column, and the split is now clean. The final model's \emph{probability} is unambiguously the best churn detector (highest AUC, Section~\ref{sec:6-quality}) -- against the Success Criteria of Section~\ref{sec:2-6}, it clears the ``AUC beats the baselines'' bar comfortably. Expected Loss has a far lower AUC by design: it trades some certainty of churn for the value of who churns, so fewer than half of its targets actually churn. But ranked by probability alone, the final model is not the best business decision: it trails the recency rule in net impact under both the reference and the realistic measurement. Ranked by \textbf{Expected Loss} -- the same model, calibrated and multiplied by each customer's value -- it meets the ``business-scenario'' criterion (positive and above the recency rule and random at the same $K$) under every measurement and every weekly budget from 1\% to 10\% with suppression. The recommended operating rule is therefore: \textbf{score churn risk with the model, rank by Expected Loss, contact a small weekly share with suppression} (Chapter~\ref{ch:deployment}). All figures remain scenario estimates under the stated assumptions (Section~\ref{sec:7-12}).

\section{Error Analysis and Limitations}
\label{sec:7-12}

Investigate false positives, false negatives, segments (one-time vs repeat, UK vs non-UK, wholesale-sized customers) and weeks where performance differs, including seasonal effects. State clearly that Online Retail II contains no campaign treatment/control information, so incremental revenue, causal retention effect and true campaign ROI cannot be identified. Business-impact results are scenario estimates.

Targeting the top 10\% by the final model's score, pooled across the test period: 57.5\% true negatives, 32.5\% false negatives, 7.6\% true positives, 2.4\% false positives -- the false-negative share is large by construction, since a 10\%-wide net cannot catch most of a $\approx$40\% base rate.

\begin{center}
\begin{tabular}{@{}p{3.6cm}p{2.6cm}p{2.6cm}@{}}
    \toprule
    Group & FP rate (among targeted) & FN rate (among non-targeted) \\
    \midrule
    One-time & 23.6\% & 54.8\% \\
    Repeat & 25.7\% & 34.6\% \\
    \midrule
    UK & 23.4\% & 36.5\% \\
    Non-UK & 36.7\% & 32.5\% \\
    \midrule
    Retail & 24.3\% & 37.0\% \\
    Wholesale-like & n/a (0 targeted) & 5.6\% \\
    \bottomrule
\end{tabular}
\end{center}

Three patterns stand out. \textbf{One-time customers} are missed far more often than repeat customers (54.8\% vs.\ 34.6\% FN rate) -- with only one purchase occasion, their feature history is thin, and the model has less signal to work with (consistent with the One-time group of Section~\ref{sec:4-6}). \textbf{Non-UK customers} see a clearly higher false-positive rate (36.7\% vs.\ 23.4\% for UK) -- the model, trained overwhelmingly on UK behavior (Section~\ref{sec:4-2}), generalizes less cleanly to the small non-UK population (9.6\% of test rows). \textbf{Wholesale-like customers} (monetary value so far at or above the train period's 99th percentile; 658 test rows) are the model's strongest case: a 5.6\% false-negative rate, far below retail's 37.0\%, and none were false positives in the targeted top 10\% -- their high, frequent purchase volume makes their churn signal unusually clear.

\textbf{Principle.} \textbf{Project-wide business-impact principle.} The model identifies and prioritizes churn risk; the retention campaign determines treatment; actual incremental revenue requires treatment/control data. Revenue at risk, revenue preserved, cost saving and net impact are simulation outputs unless supported by observed campaign outcomes.
